\documentclass{article}
\usepackage{graphicx} % Required for inserting images
\usepackage{amsmath}
\title{Rocketry equation}
\author{주안 김}
\date{July 2026}

\begin{document}

\maketitle

\section{rocketry equation}


\[
\textbf{Nominal Operating Point}
\]

Let

\[
\mathbf{x}_0=(W_0,V_0,\Delta T_0)
\]

denote the nominal launch conditions, and let

\[
\varepsilon\in\mathbb{R},
\qquad
|\varepsilon|\ll1
\]

represent a small perturbation from the nominal state.

Assume

\[
W(\varepsilon)
=
W_0
\left(
1+\sum_{n=1}^{\infty}\alpha_n\varepsilon^n
\right),
\]

\[
V(\varepsilon)
=
V_0
+
\sum_{n=1}^{\infty}\beta_n\varepsilon^n,
\]

\[
\Delta T(\varepsilon)
=
\Delta T_0
+
\sum_{n=1}^{\infty}\gamma_n\varepsilon^n.
\]

The apogee model is

\[
H(W,V,\Delta T)
=
A
+
B(\Delta T)
+
C\ln(W)
+
D\cos(V).
\]

Substituting the perturbation expansions,

\[
\begin{aligned}
H(\varepsilon)
=
&
A
+
B
\left(
\Delta T_0
+
\sum_{n=1}^{\infty}
\gamma_n\varepsilon^n
\right)
\\
&
+
C
\ln
\left[
W_0
\left(
1+\sum_{n=1}^{\infty}
\alpha_n\varepsilon^n
\right)
\right]
\\
&
+
D
\cos
\left(
V_0
+
\sum_{n=1}^{\infty}
\beta_n\varepsilon^n
\right).
\end{aligned}
\]

Using

\[
\ln(ab)=\ln(a)+\ln(b),
\]

gives

\[
\ln(W)
=
\ln(W_0)
+
\ln
\left(
1+\sum_{n=1}^{\infty}
\alpha_n\varepsilon^n
\right).
\]

Applying the Taylor series

\[
\ln(1+x)
=
\sum_{k=1}^{\infty}
(-1)^{k+1}
\frac{x^k}{k},
\qquad
|x|<1,
\]

where

\[
x
=
\sum_{n=1}^{\infty}
\alpha_n\varepsilon^n,
\]

yields

\[
\ln(W)
=
\ln(W_0)
+
\sum_{k=1}^{\infty}
(-1)^{k+1}
\frac{1}{k}
\left(
\sum_{n=1}^{\infty}
\alpha_n\varepsilon^n
\right)^k.
\]

Similarly, expanding about the nominal wind speed \(V_0\),

\[
\begin{aligned}
\cos(V)
=
&
\cos(V_0)
-
\sin(V_0)
\left(
\sum_{n=1}^{\infty}
\beta_n\varepsilon^n
\right)
\\
&
-
\frac{\cos(V_0)}{2!}
\left(
\sum_{n=1}^{\infty}
\beta_n\varepsilon^n
\right)^2
\\
&
+
\frac{\sin(V_0)}{3!}
\left(
\sum_{n=1}^{\infty}
\beta_n\varepsilon^n
\right)^3
+\cdots.
\end{aligned}
\]

Hence,

\[
\begin{aligned}
H(\varepsilon)
=
&
A
+
B\Delta T_0
+
C\ln(W_0)
+
D\cos(V_0)
\\
&
+
\sum_{n=1}^{\infty}
K_n
\varepsilon^n,
\end{aligned}
\]

where

\[
K_n
=
B\gamma_n
+
CL_n
+
DM_n,
\]

with

\[
L_n
\]

the coefficient of
\(\varepsilon^n\)
from the logarithmic expansion, and

\[
M_n
\]

the coefficient of
\(\varepsilon^n\)
from the cosine expansion.

Finally,

\[
\boxed{
H(\varepsilon)
=
K_0
+
\sum_{n=1}^{\infty}
K_n
\varepsilon^n
}
\]

where

\[
K_0
=
A
+
B\Delta T_0
+
C\ln(W_0)
+
D\cos(V_0).
\]

\[
\textbf{Local Sensitivity with Respect to Weight}
\]

The apogee model is

\[
H(W,V,\Delta T)
=
A
+
B(\Delta T)
+
C\ln(W)
+
D\cos(V).
\]

The first partial derivative with respect to weight is

\[
\frac{\partial H}{\partial W}
=
\frac{C}{W}.
\]

Hence, the local sensitivity of the apogee to weight is inversely
proportional to the rocket weight. In particular,

\[
\lim_{W\to0^+}
\frac{\partial H}{\partial W}
=\infty,
\]

indicating that small variations in lighter rockets produce
significantly larger changes in the predicted apogee.

\vspace{0.5cm}

\[
\textbf{Local Curvature}
\]

Differentiating a second time,

\[
\frac{\partial^2H}{\partial W^2}
=
-\frac{C}{W^2}.
\]

Since

\[
\frac{\partial^2H}{\partial W^2}<0
\qquad (C>0),
\]

the logarithmic contribution is locally concave, implying diminishing
returns as the rocket weight increases.

Furthermore,

\[
\left|
\frac{\partial^2H}{\partial W^2}
\right|
=
\frac{|C|}{W^2},
\]

which decreases quadratically with increasing weight. Thus, the model
is most strongly curved in the low-weight regime and gradually becomes
locally linear as \(W\) increases.

\vspace{0.5cm}

Consequently, the first-order Taylor approximation is expected to be
most accurate for heavier rockets, while higher-order terms become
increasingly important for lighter rockets due to the larger local
curvature.

\[
\textbf{Launch Manifold}
\]

\[
M=\{(W,V,\Delta T)\}, \qquad
H:M\rightarrow\mathbb{R},
\]

\[
H(W,V,\Delta T)
=
A+B\Delta T+C\ln(W)+D\cos(V).
\]

\[
\textbf{Coordinate Map}
\]

\[
u=\ln(W),
\qquad
du=\frac{dW}{W}.
\]

Hence

\[
H(u,V,\Delta T)
=
A+B\Delta T+Cu+D\cos(V).
\]

\[
\textbf{Metric}
\]

\[
g
=
\frac{dW^2}{W^2}
+
\frac{dV^2}{V_0^2}
+
\frac{d(\Delta T)^2}{T_0^2}.
\]

Equivalently,

\[
g
=
du^2
+
\frac{dV^2}{V_0^2}
+
\frac{d(\Delta T)^2}{T_0^2}.
\]

\[
\textbf{Christoffel Symbol}
\]

\[
g_{WW}
=
\frac1{W^2},
\]

\[
\Gamma^W_{WW}
=
\frac12g^{WW}
\left(
2\partial_Wg_{WW}
-
\partial_Wg_{WW}
\right)
=
-\frac1W.
\]

\[
\textbf{Geodesic Equation}
\]

\[
\frac{d^2x^k}{ds^2}
+
\Gamma^k_{ij}
\frac{dx^i}{ds}
\frac{dx^j}{ds}
=
0.
\]

For the weight coordinate,

\[
\frac{d^2W}{ds^2}
-
\frac1W
\left(
\frac{dW}{ds}
\right)^2
=
0.
\]

\[
\textbf{Covariant Hessian}
\]

\[
\nabla_i\nabla_jH
=
\partial_i\partial_jH
-
\Gamma^k_{ij}\partial_kH.
\]

\[
\frac{\partial H}{\partial W}
=
\frac{C}{W},
\]

\[
\frac{\partial^2H}{\partial W^2}
=
-\frac{C}{W^2},
\]

\[
\frac{\partial^2H}{\partial V^2}
=
-D\cos(V),
\]

\[
\frac{\partial^2H}{\partial(\Delta T)^2}
=
0.
\]

\[
\textbf{Covariant Expansion}
\]

\[
H(p)
=
H(p_0)
+
\nabla_iH\,\xi^i
+
\frac12
\nabla_i\nabla_jH
\,\xi^i\xi^j
+
O(\|\xi\|^3).
\]

\section*{application}
\[
\textbf{Apogee Sensitivity Model}
\]

\[
H=H(W,V,\Delta T)
\]

\[
\Delta H
\approx
\nabla H\cdot\Delta x
+
\frac12
\Delta x^T
(\nabla^2H)
\Delta x
\]

where

\[
\Delta x=
\begin{bmatrix}
\Delta W\\
\Delta V\\
\Delta T
\end{bmatrix}.
\]

\[
\textbf{Sensitivity}
\]

\[
\nabla H
=
\begin{bmatrix}
\dfrac{\partial H}{\partial W}\\[6pt]
\dfrac{\partial H}{\partial V}\\[6pt]
\dfrac{\partial H}{\partial \Delta T}
\end{bmatrix}
\]

\[
S_i=
\left|
\frac{\partial H}{\partial x_i}
\right|
\]

\[
\eta_i=
\frac{S_i}{\sum_jS_j}
\]

\[
\textbf{Uncertainty Propagation}
\]

\[
\sigma_H^2
=
\nabla H^T
\Sigma
\nabla H
\]

where

\[
\Sigma=
\begin{bmatrix}
\sigma_W^2&0&0\\
0&\sigma_V^2&0\\
0&0&\sigma_T^2
\end{bmatrix}.
\]

\[
\textbf{Optimization}
\]

\[
\max H(W,V,\Delta T)
\]

\[
\nabla H=0
\]

\[
\nabla^2H
\begin{cases}
<0,&\text{maximum}\\
>0,&\text{minimum}
\end{cases}
\]

\[
\textbf{Launch Condition Manifold}
\]

\[
M=\{(W,V,\Delta T)\}
\]

\[
H:M\rightarrow\mathbb{R}
\]

\[
ds^2
=
\frac{dW^2}{W^2}
+
\frac{dV^2}{V_0^2}
+
\frac{d(\Delta T)^2}{T_0^2}
\]

\[
\textbf{Geodesics}
\]

\[
\frac{d^2x^k}{ds^2}
+
\Gamma^k_{ij}
\frac{dx^i}{ds}
\frac{dx^j}{ds}
=0
\]

\[
\textbf{Curvature}
\]

\[
\nabla_i\nabla_jH
=
\partial_i\partial_jH
-
\Gamma^k_{ij}\partial_kH
\]

\[
\frac{\partial^2H}{\partial W^2}
=
-\frac{C}{W^2}
\]

\[
\frac{\partial^2H}{\partial V^2}
=
-D\cos(V)
\]

\[
\textbf{Performance Metric}
\]

\[
J(W,V,\Delta T)
=
\frac{H(W,V,\Delta T)}{W}
\]

\[
\max J
\]

\[
\nabla J=0
\]

\section*{Solve}

\[
\textbf{1. Data Collection}
\]

Given experimental/simulation data:

\[
(W_i,V_i,\Delta T_i,H_i),
\qquad i=1,2,\dots,n
\]

where

\[
H_i=A+B\Delta T_i+C\ln(W_i)+D\cos(V_i).
\]

---

\[
\textbf{2. Matrix Form}
\]

Define

\[
X=
\begin{bmatrix}
1&\Delta T_1&\ln(W_1)&\cos(V_1)\\
1&\Delta T_2&\ln(W_2)&\cos(V_2)\\
\vdots&\vdots&\vdots&\vdots\\
1&\Delta T_n&\ln(W_n)&\cos(V_n)
\end{bmatrix}
\]

\[
\theta=
\begin{bmatrix}
A\\B\\C\\D
\end{bmatrix}
\]

\[
Y=
\begin{bmatrix}
H_1\\H_2\\\vdots\\H_n
\end{bmatrix}.
\]

Therefore,

\[
X\theta=Y.
\]

---

\[
\textbf{3. Least Squares Solution}
\]

For

\[
n>4,
\]

solve:

\[
\boxed{
\theta=(X^TX)^{-1}X^TY
}
\]

giving

\[
\boxed{
(A,B,C,D)
}
\]

---

\[
\textbf{4. Construct Model}
\]

\[
\boxed{
H(W,V,\Delta T)
=
A+B\Delta T+C\ln(W)+D\cos(V)
}
\]

---

\[
\textbf{5. Sensitivity}
\]

\[
\nabla H
=
\begin{bmatrix}
\dfrac{\partial H}{\partial W}\\
\dfrac{\partial H}{\partial V}\\
\dfrac{\partial H}{\partial \Delta T}
\end{bmatrix}
\]

\[
\boxed{
\nabla H
=
\begin{bmatrix}
\dfrac{C}{W}\\
-D\sin(V)\\
B
\end{bmatrix}
}
\]

---

\[
\textbf{6. Curvature}
\]

\[
\nabla^2H
=
\begin{bmatrix}
-\dfrac{C}{W^2}&0&0\\
0&-D\cos(V)&0\\
0&0&0
\end{bmatrix}
\]

---

\[
\textbf{7. Taylor Approximation}
\]

Around

\[
x_0=(W_0,V_0,\Delta T_0)
\]

\[
\boxed{
H(x)
\approx
H(x_0)
+
\nabla H(x_0)\Delta x
+
\frac12
\Delta x^T
\nabla^2H(x_0)
\Delta x
}
\]

---

\[
\textbf{8. Optimization}
\]

Find

\[
\nabla H=0
\]

and classify using

\[
\nabla^2H.
\]

\[
\boxed{
\text{Optimal launch condition}
}
\]


\[
\textbf{Configuration Manifold}
\]

\[
M=(0,\infty)\times\mathbb{R}\times\mathbb{R}
\]

\[
H:M\rightarrow\mathbb{R}
\]

\[
H(W,V,\Delta T)
=
A+B\Delta T+C\ln(W)+D\cos(V)
\]

---

\[
\textbf{Level Sets}
\]

\[
\Sigma_c
=
H^{-1}(c)
=
\left\{
x\in M:
H(x)=c
\right\}
\]

\[
\Sigma_c
=
\left\{
(W,V,\Delta T):
A+B\Delta T+C\ln(W)+D\cos(V)=c
\right\}
\]

---

\[
\textbf{Regular Values}
\]

\[
\nabla H\neq0
\]

\[
\Longrightarrow
\]

\[
\Sigma_c
\text{ is a smooth }2\text{-manifold.}
\]

---

\[
\textbf{Critical Points}
\]

\[
\nabla H=0
\]

\[
\begin{cases}
\dfrac{C}{W}=0\\[6pt]
-D\sin(V)=0\\[6pt]
B=0
\end{cases}
\]

---

\[
\textbf{Hessian}
\]

\[
\nabla^2H
=
\begin{pmatrix}
-\dfrac{C}{W^2} &0&0\\
0&-D\cos(V)&0\\
0&0&0
\end{pmatrix}
\]

\[
\det(\nabla^2H)\neq0
\]

\[
\Longrightarrow
\]

\[
H
\text{ is Morse (locally).}
\]

---

\[
\textbf{Sublevel Sets}
\]

\[
M_c
=
\{x\in M:H(x)\le c\}
\]

\[
M_{c_1}
\subseteq
M_{c_2},
\qquad
c_1<c_2
\]

---

\[
\textbf{Boundary}
\]

\[
\partial M_c
=
\Sigma_c
=
H^{-1}(c)
\]

---

\[
\textbf{Betti Numbers}
\]

\[
\beta_0
=
\#(\text{connected components})
\]

\[
\beta_1
=
\#(\text{independent loops})
\]

\[
\beta_2
=
\#(\text{enclosed voids})
\]

---

\[
\textbf{Critical Values}
\]

\[
c^\ast
=
H(x^\ast),
\qquad
\nabla H(x^\ast)=0
\]

\[
c=c^\ast
\Longrightarrow
\text{Topology Change}
\]

---

\[
\textbf{Homotopy}
\]

\[
\Sigma_{c_1}
\simeq
\Sigma_{c_2},
\qquad
c_1,c_2
\notin
\{c^\ast\}
\]

---

\[
\textbf{Persistent Filtration}
\]

\[
\varnothing
=
M_{c_0}
\subseteq
M_{c_1}
\subseteq
\cdots
\subseteq
M_{c_n}
=
M
\]

\[
\beta_k(c)
=
\operatorname{rank}
H_k(M_c)
\]

---

\[
\boxed{
\begin{aligned}
M
&\xrightarrow{\;H\;}
\mathbb{R}
\\
&
\Downarrow
\\
H^{-1}(c)
&\Longrightarrow
\Sigma_c
\\
&
\Downarrow
\\
\nabla H=0
&\Longrightarrow
c^\ast
\\
&
\Downarrow
\\
\text{Topology Transition}
\end{aligned}
}
\]


\[
\textbf{Time-Dependent Variables}
\]

\[
W=W(t),
\qquad
V=V(t),
\qquad
\Delta T=\Delta T(t)
\]

\[
H(t)
=
A+B\Delta T(t)+C\ln(W(t))+D\cos(V(t))
\]

---

\[
\textbf{Chain Rule}
\]

\[
\frac{dH}{dt}
=
\frac{\partial H}{\partial W}\dot W
+
\frac{\partial H}{\partial V}\dot V
+
\frac{\partial H}{\partial(\Delta T)}\dot T
\]

\[
\boxed{
\frac{dH}{dt}
=
\frac{C}{W}\dot W
-
D\sin(V)\dot V
+
B\dot T
}
\]

---

\[
\textbf{Second Derivative}
\]

\[
\boxed{
\frac{d^2H}{dt^2}
=
-\frac{C}{W^2}\dot W^2
+
\frac{C}{W}\ddot W
-
D\cos(V)\dot V^2
-
D\sin(V)\ddot V
+
B\ddot T
}
\]

---

\[
\textbf{Differential Operator}
\]

\[
L=\frac{d}{dt}
\]

\[
LH
=
\frac{C}{W}\dot W
-
D\sin(V)\dot V
+
B\dot T
\]

\[
L^2H
=
\frac{d^2H}{dt^2}
\]

---

\[
\textbf{Vector Field}
\]

\[
X
=
\dot W\frac{\partial}{\partial W}
+
\dot V\frac{\partial}{\partial V}
+
\dot T\frac{\partial}{\partial(\Delta T)}
\]

\[
\boxed{
\mathcal L_XH
=
X(H)
=
\frac{dH}{dt}
}
\]

---

\[
\textbf{Exponential Weight}
\]

\[
\dot W=-kW
\]

\[
W(t)=W_0e^{-kt}
\]

\[
\frac{C}{W}\dot W
=
-kC
\]

---

\[
\textbf{Linear Temperature}
\]

\[
\Delta T
=
\Delta T_0+\alpha t
\]

\[
\dot T=\alpha
\]

\[
B\dot T
=
B\alpha
\]

---

\[
\textbf{Oscillatory Wind}
\]

\[
V
=
V_0+\beta\sin(\omega t)
\]

\[
\dot V
=
\beta\omega\cos(\omega t)
\]

\[
-D\sin(V)\dot V
=
-D\beta\omega
\sin(V)
\cos(\omega t)
\]

---

\[
\textbf{Fourier Extension}
\]

\[
F(V)
=
a_0
+
\sum_{n=1}^{\infty}
\left(
a_n\cos(nV)
+
b_n\sin(nV)
\right)
\]

\[
H
=
A
+
B\Delta T
+
C\ln(W)
+
F(V)
\]

---

\[
\textbf{Fourier Coefficients}
\]

\[
a_n
=
\frac1\pi
\int_{-\pi}^{\pi}
F(V)\cos(nV)\,dV
\]

\[
b_n
=
\frac1\pi
\int_{-\pi}^{\pi}
F(V)\sin(nV)\,dV
\]

---

\[
\textbf{Generalized Model}
\]

\[
H
=
A
+
B\Delta T
+
C\ln(W)
+
\sum_{n=1}^{\infty}
\left(
a_n\cos(nV)
+
b_n\sin(nV)
\right)
\]

---

\[
\textbf{Generalized Dynamics}
\]

\[
\boxed{
\frac{dH}{dt}
=
\frac{C}{W}\dot W
+
B\dot T
+
\sum_{n=1}^{\infty}
n
\left(
-a_n\sin(nV)
+
b_n\cos(nV)
\right)
\dot V
}
\]

---

\[
\textbf{Asymptotic Expansion}
\]

\[
H
=
H_0
+
\varepsilon H_1
+
\varepsilon^2H_2
+
O(\varepsilon^3)
\]

---

\[
\boxed{
H
\rightarrow
\frac{dH}{dt}
\rightarrow
\frac{d^2H}{dt^2}
\rightarrow
\mathcal L_XH
\rightarrow
F(V)
\rightarrow
\text{Fourier Dynamics}
}
\]
\end{document}
